\section{Trade Theory}




\begin{frame}{Why Trade?}
    \mytext{
    \begin{itemize}
        \item Nobody has to buy, and
        \item nobody has to sell.
    \end{itemize}

    I could grow my own food, make my own shoes... but obviously, trading makes sense.

    The wealth of nation is based on the division of labour (Adam \noun{Smith}, 1776).

    Typically, we sell just some products (or services), but we buy a myriad of things and services.

    This is true for individuals, companies, and even countries.
    }
\end{frame}

\begin{frame}{Why Trade?}
    \mytext{
    \begin{itemize}
        \item goods we otherwise would not have
        \begin{itemize}
            \item different climate
            \item different resources
            \item different know-how
        \end{itemize}
        \item goods we have too much of
        \begin{itemize}
            \item dumping, overcapacities
        \end{itemize}
        \item goods we sell to buy other goods
        \begin{itemize}
            \item produce for the market
            \item to exchange for other things
            \item (profit) maximisation motive
        \end{itemize}
    \end{itemize}
    }
\end{frame}

\begin{frame}{Why Trade?}
    \mycomment{
    We do not just trade things that we happen to have, but we have things so that we can trade. The decision what to trade and what to produce are  intrinsically linked.
    }
    \mytext{
    \begin{itemize}
        \item to resell the goods: profits / artbitrage
        \item to complement our demand
        \item to consume
        \item to use to manufacture
        \begin{itemize}
            \item raw materials
            \item energy
            \item investment
        \end{itemize}
    \end{itemize}
    }
\end{frame}
\begin{frame}{Why Trade?}
    \mycomment{
    Our decisions what to produce - for a market, that is, to trade - depends on our endowment. What can we produce, and what can we produce best?

    The production decision will depend on our internal opportunity costs - what we can do best; but will also depend on the relative price of the goods we produce, that is: relative to the price relationships on the market.

    And these market prices again will be an outcome of the markets, on which many parties with different endowment interact.
    
    }
    \mytext{
    Production Factors: 
    
    \begin{itemize}
        \item goods / resources
        \item capital
        \item labour
    \end{itemize}
    }
\end{frame}
\begin{frame}{Why Trade?}
    \mycomment{
        Trading would not take place if it was not for the trading parties benefit. After all, trading is not robbery. Nobody is forced to buy or sell.

        That being said, people may in fact be in the situation where they have to trade to survive. If I have no water, I must sell something to get it, or I will die. So there may be situations where people actually are desperate to sell.

        In economics, desperate situations are corner conditions. And these are often sub-optimal. It would be better if nobody be forced to act in a certain way, but anyone could participate in some kind of optimisation process.
    }
    \mytext{
    \begin{itemize}
        \item Someone offers something cheaper
        \item Someone (abroad) is willing to pay more 
        \item Someone (abroad) is willing to pay as much
        \item We can all profit
        \item No man is an island
    \end{itemize}
    }
\end{frame}


\begin{frame}{Thünen's Rings}
    \showmytikz{../shared/tikz/tikz-thuenen}
    \mycomment{
        Johann Heinrich von \noun{Thünen} observed 1842 that as you travel outbound from a city, the land is used for different purposes. This forms the basis for many models of Economic Geography.
    }
\end{frame}
\begin{frame}{Thünen's Rings}
    \showmytikz{../shared/tikz/tikz-thuenen-grafisch}
    \mycomment{
        \noun{Thünen} explained this by the role of transportation costs, which were high in his time (carriage with oxen). 

        However, today we have multiple modes of transport. Shipping is cheap, but also relatively slow. Fresh goods like flowers or fresh fruit are today often transported via air fright from different climate zones.        

        However, the model still serves well to explain e.g. housing rents in the vicinity of cities.
    }
\end{frame}


\begin{frame}{Gravity Model of Trade}
    \showmytikz{../shared/tikz/tikz-gravitymodel}
    
    \mycomment{
    Trade between countries depends on...
    \begin{itemize}
        \item size of the countries
        \item distance of the countries
    \end{itemize}

    \begin{equation}
        T_y=A\times Y_i \times \frac{Y_j}{D_{ij}}
    \end{equation}

    \begin{description}
        \item[A] const.
        \item[Y] GDP (of $i$ or $j$ )
        \item[D] Distance (between $i$ and $j$)
    \end{description}

    Distance would be measured in transportation / transaction costs, including risk, trade barriers etc.
  
    }
\end{frame}



\begin{frame}{Trade Openness}
\mycomment{
    Trade openness measures how large a percentage of gross domestic product is traded across the country's border, that is, trade volume is set in relation to GDP. 

    Alternatively, sometimes the export ratio or the import ratio is given. These can, of course, be then put into relation to one another, showing imbalances in the trade balance.
}
\mytext{

    \begin{equation}
        \text{Openness} = \frac{\text{Exp}+\text{Imp}}{\text{GDP}}
    \end{equation}  
}
\end{frame}
\begin{frame}{Trade Openness}
\mycomment{
    Larger economies tend to be "less open", but this is a denominator-effect. Their (large) trade is small compared to their large gross domestic product. Smaller economies, however, have typically a larger trade openness. This also means that they are more affected by external shocks.
}
\mytable{
    \begin{tabular}{lrl}
        \toprule
         World                &  62 & \myblock[0.2ex]{62}\\
         EU                   & 103 & \myblock[0.2ex]{103}\\
         \midrule
         Brasil               &  38 & \myblock[0.2ex]{38}\\
         Chile                &  75 & \myblock[0.2ex]{75}\\
         China                &  37 & \myblock[0.2ex]{37}\\
         Congo Rep.           &  98 & \myblock[0.2ex]{98}\\
         Germany              &  88 & \myblock[0.2ex]{88}\\
         High Income          &  67 & \myblock[0.2ex]{67}\\
         Hong Kong            & 352 & \\
         Indonesia            &  45 & \myblock[0.2ex]{45}\\
         Least developed (UN) &  51 & \myblock[0.2ex]{51}\\
         Mexico               &  88 & \myblock[0.2ex]{88}\\
         UK                   &  69 & \myblock[0.2ex]{69}\\
         USA                  &  27 & \myblock[0.2ex]{27}\\
         \bottomrule
        \bottomrule
    \end{tabular}
    
    Data for 2022 \cite{worldbankdatabank}
}
\end{frame}

\begin{frame}{Exchange: Edgeworth Box}
    \showmytikz{../shared/tikz/tikz-edgeworth}
    \mycomment{
    We can plot the preferences of $A$ so that more of $X,Y$ is always better (assuming $X,Y$ to be economically normal goods). 

    The production point $P$ for $A$ will depend on $A$'s preferences - on how much they like to have the goods, and also how much effort they like to spend having them.
    }
\end{frame}
\begin{frame}{Exchange: Edgeworth Box}
    \showmytikz[1,2]{../shared/tikz/tikz-edgeworth}
    \mycomment{
    We can do the same for  $B$ and see that $A$ would offer $B$ some of his $Y$ in exchange for $X$ to reach a higher level of utility. 

    They could agree on point $Q$, where $A$ has no real advantage, but is just as well off as before. But at point $R$ they will be better off. So in a symmetrical case, there are many Pareto optimal outcomes.


    }
\end{frame}
\begin{frame}{Exchange: Edgeworth Box}
    \showmytikz[1,2,3,4]{../shared/tikz/tikz-edgeworth}
    \mycomment{


    The blue line is called the exchange curve.

    Any point on the trading curve is stable, as any movement would necessarily render one of the two parties worse off than before. 

    This can be used to demonstrate the First Fundamental Theorem of Welfare Economics: that free exchange is Pareto-optimal.

    This does, however, only demonstrate that trade, meaning exchange, is a good thing. But how does this influence production decisions, that is: when we do not just exchange stuff we happen to have, but when we produce for the market?
    
    }
\end{frame}




\begin{frame}{Differences Lead to Trade}
    \showmytikz[1]{../shared/tikz/tikz-thuenen-twice}
    \mycomment{
    If two places be exactly identical, there would be not much incentive for trade. Maybe the farmers would exchange one cow for another for reasons of variety.

    Note that "exactly the same" here implies that not only areas of and means for production are identical, but also the number of people, and hence local wages.
    }
\end{frame}
\begin{frame}{Differences Lead to Trade}
    \showmytikz[1,2,3]{../shared/tikz/tikz-thuenen-twice}
    \mycomment{
    If, however, a potential trading partner is different, the situation changes. 
    
    In our picture, the third potential trading partner (island) has no agriculture - maybe due to rocky terrain. Surely the people living there would be happy to trade some of their livestock, wood or garden produce for wheat, bread, or rice!
    }
\end{frame}

\begin{frame}{Differences Lead to Trade}
    \mytext{
    Differences lead to trade:
    \begin{itemize}
        \item absolute difference: swap
        \item production advantages:
        \begin{itemize}
            \item absolute advantage
            \item comparative advantage
        \end{itemize}
    \end{itemize}
    }
    \mycomment{
        Not only an absolute difference can lead to trade. Trade can make sense and will take place even if both countries (or even persons who trade) can produce the same products. It is sufficient if one party is better at something than the other.

        There is absolute advantage, where $A$ has a higher productivity in the production of $X$ than $B$:

        \begin{equation}
            \Pi_{A}(X) > \Pi_B(X) \quad \lor \quad \Pi_{B}(Y) > \Pi_B(Y) 
        \end{equation}

        And there is relative advantage, where the ratios of the productivities in the production of $X,Y$ between $A,B$ are different:

        \begin{equation}
            \frac{\Pi_A(X)}{\Pi_A(Y)}>\frac{\Pi_B(X)}{\Pi_B(Y)}            
        \end{equation}
    }
\end{frame}


\begin{frame}{Absolute Advantage: Visually}
    \showmytikz[1]{../shared/tikz/tikz-exchangesmallandsmall}
    \mycomment{
    The continent has more fruit than fish, trading internally at 4:1. So you pay a lot of fruit for one fish.

    The small island has more fish than fruit. Fruit and fish here trade at a rate of 1:4. So on the island, you pay a lot of fish for one fruit.

    Now, if the island were to trade with the continent, every island fisher will be better off by offering their catch on the continent.

    On the continent, they get four fruits for just one fish! That is 16 times as much as on the island! 

    The situation is similar for the continental farmers. 
    }
\end{frame}
\begin{frame}{Absolute Advantage: Visually}
    \showmytikz[1,2]{../shared/tikz/tikz-exchangesmallandsmall}
    \mycomment{
    \begin{itemize}
        \item The fruit producer on the island will be outpriced by the continental competition.  The same applies for fishermen on the continent.
        \item In the end, specialisation on both island and continent will be total.
        \item If there are transport costs, these will diminish the effect; but as long as the arbitrage is sufficient, the transport will be paid for.
        \item The island fishermen will compete with the fishermen on the continent; but the fishmongers will work with both, as they are in the same industry. 
        \item Likewise, the former farmers on the island will now probably sell imported fruit.
        \item The traders will ship fish in one direction and fruit in the other.
        \item The new trading rate of fish to fruit is now closer to one, resulting in a healthier diet and larger utility in both industries and both countries.
    \end{itemize}
    }
\end{frame}
\begin{frame}{Absolute Advantage: Visually}
    \showmytikz[1,2,3]{../shared/tikz/tikz-exchangesmallandsmall}
     \mycomment{
    \begin{itemize}
        \item The island farmers will be better off to also start fishing. Island fishing will go up.
        \item There will be more fish on offer on the continent. As a result, the price of fish (relative to fruit) will  fall (and the price of fruit rise, relatively). 
        \item This will have the fishermen on the continent go into the fruit industry, which is now  more lucrative.
        \item In the end, there will be a new equilibrium both on the continent and on the island, with more fish and more fruit on both.  
        \item Fruit producers on the island and fishermen on the continent may (at first) oppose trade. Fishermen on the island and fruit producers on the continent will be in favour.
        \item Both will lose an production industry and become dependent on each other. But both will also have a trading industry, and have both fish and fruit.
    \end{itemize}
    }
\end{frame}


\begin{frame}{Absolute Advantage: II}
    \showmytikz[1]{../shared/tikz/tikz-absoluteadvantage}
    \mycomment{
        Without trade, Red will produce and consume a combination of $X,Y$ at point $A$.

        For example, Red produces more fish ($Y$) and less bread ($X$). That is because for Red it is easy to fish and hard to grow grain.
    }
\end{frame}

\begin{frame}{Absolute Advantage: II}
    \showmytikz[2]{../shared/tikz/tikz-absoluteadvantage}
    \mycomment{
       Without trade, Blue will produce and consume a combination of $X,Y$ at point $B$; so more bread ($X$) than fish ($Y$).
       
       For Blue producing fish ($Y$) is more expensive in absolute terms than producing bread ($X$).  
    }
\end{frame}

\begin{frame}{Absolute  Advantage: II}
    \showmytikz[1,2]{../shared/tikz/tikz-absoluteadvantage}
    \mycomment{

        We see that Red has an absolute advantage in the production of $Y$, and Blue in the production of $X$ relative to one another.
        
        If they trade, each will be happy to offer the respective good at a price the other one cannot hold. 

        Consequentially, they will each focus on whatever it is they can do best.
    }
\end{frame}

\begin{frame}{Absolute Advantage: II}
    \showmytikz[1,2,3]{../shared/tikz/tikz-absoluteadvantage}
    \mycomment{
        If Red and Blue trade, they will enjoy more of $X, Y$  than they could without trade.

        They unite or pool their production frontiers: their coordination results in combined optimisation.

        As a result, each will focus on producing what he/she can do best, that is: cheaper; as he/she can exchange the other good for this. Red will be a fisher, and Blue a farmer.

        In the new equilibrium, both will enjoy a new mix of $X,Y$ that makes them both better off. Red will now eat less fish than before, but more bread, and Blue will enjoy more fish and less bread.

        Total production is likely, but not necessarily, higher than before. But the utility reached under the new distribution will be higher than before. It could also be that both eat better, i.e. a more diverse diet, with less work.
    }
\end{frame}


\begin{frame}{Comparative Advantage: Visually}
    \showmytikz[1]{../shared/tikz/tikz-exchangelargeandsmall}
    \mycomment{
    Trade and specialisation are a good thing, even if one of the trading countries has an absolute advantage in both industries.

    The rich and large continent has much fish and much fruit; they trade at a rate of 1:1.

    The poor and small island has some fish, but little fruit. Fruit and fish here trade at a rate of 1:4.

    (The numbers here are: 16 fruit and 16 fish on the continent; 1 fruit on the island; 4 fish on the island.)

    Now, if the island were to trade with the continent, every island fisher will still be better off by offering their catch on the continent, because they will get four times as much fruit for it than on the island.
    }
\end{frame}
\begin{frame}{Comparative Advantage: Visually}
    \showmytikz[1,2]{../shared/tikz/tikz-exchangelargeandsmall}
    \mycomment{
    \begin{itemize}
        \item The fruit producer on the island will (again) be outpriced by the continental fruit producers. 
        \item Specialisation on the island will (again) be total, but specialisation on the continent only partial.
        \item If there are transport costs, these will diminish the effect; but as long as the arbitrage is sufficient, the transport will be paid for.
        \item The island fishermen will compete with the fishermen on the continent; but the fishmongers will work with both, as they are in the same industry.
        \item The traders will ship both fish and fruit.
        \item The new trading rate of fish to fruit on the island is now also closer to one, resulting in a healthier diet and larger utility for the islanders.
        \item If the island is very small, the impact on the continent is negligible. 
    \end{itemize}
    }
\end{frame}
\begin{frame}{Comparative Advantage: Visually}
    \showmytikz[1,2,3]{../shared/tikz/tikz-exchangelargeandsmall}
     \mycomment{
    \begin{itemize}
        \item The island farmers will be better off to also start fishing. Island fishing will go up.
        \item There will be more fish on offer on the continent (even though not very much more). As a result, the price of fish (relative to fruit) will slightly fall (and the price of fruit rise, relatively). 
        \item This will have some fishermen on the continent go into the fruit industry, which is now (slightly) more lucrative.
        \item In the end, there will be a new equilibrium both on the continent and on the island, with more fish and more fruit on both. The change on the island, however, is more dramatic.
        \item Fruit producers on the island and fishermen on the continent may oppose trade. Fishermen on the island and fruit producers on the continent will be in favour.
        \item Island will lose its fruit industry and become somehow dependent on trade and fruit import.
    \end{itemize}
    }
\end{frame}

\begin{frame}{Comparative Advantage}
    \showmytikz[1,2]{../shared/tikz/tikz-comparativeadvantage}
    \mycomment{
    Now consider the case where blue is just as good at producing $X$ as red, but much worse at producing $Y$.

    Still, trade would be good. If Blue and Red combine their efforts so that Blue does only what they can do best, the trade can enable Red to then do more of what Red can do best.

    The fact that comparative, not absolute, advantage and hence opportuntiy costs of the trading partners alone stipulates trade was first observed by David \noun{Ricardo}.
    }
\end{frame}
\begin{frame}{Comparative Advantage}
    \showmytikz[1,2,3]{../shared/tikz/tikz-comparativeadvantage}
    \mycomment{
    When Red and Blue trade, they combine their production frontiers (the new, thick line). 

    Blue will now only do what he can do best and only produce $X$ (at $C$), much of which (but not all) he will trade with Red. 

    Red may only do what he can do best - but maybe he would like to have even more $X$ than Blue can provide, so he may still produce some $X$ (at $D$). So the less productive one will specialise totally, while the more productive one will likely still engage in both activities.

    Note that both will have more of each $X,Y$ than before - and a more balanced diet.

    This shows that with trade, \textbf{everyone should focus on what they can do best.}

    This can lead to the situation where you focus on something (like raw materials) forever.
    
    }
\end{frame}

\begin{frame}{Comparative Advantage}
    \mytext{
    \begin{itemize}
        \item Under trade, specialisation on what one can do best is a welfare maximising strategy for all involved - at least short term.
        \item A smaller and/or less efficient trading partner is likely to specialise in full
        \item The larger and/or more efficient trading partner will likely only specialise partly
    \end{itemize}
    }
\end{frame}
\begin{frame}{Comparative Advantage}
    \mytext{
    \begin{itemize}
        \item The exchange rate (terms of trade) of goods will likely be determined by (and resemble) the internal rate of substitution (opportunity costs) of the larger, more powerful partner
        \item long-term effects such as technological progress (increase in efficiency) may change the situation in either way
        \item relative price changes anywhere will affect the outcomes
        \item specialisation creates dependency and increases risk exposure
    \end{itemize}
    }
\end{frame}



\begin{frame}{Evidence?}
\mytext{
    The theory of comparative advantage, since it was first stated by David \noun{Ricardo}, has been a prominent argument for free trade---and indeed for colonialism. There are, however, signs that the theory does not hold, or at least does not suffice to explain foreign trade.
    \begin{itemize}
        \item Leontieff paradox: 
        \item[] there is more trade between \textbf{similar} economies 
        \item Linder hypothesis to explain this: 
        \item[] trade is demand based
        \item What about global value chains (GVC)?
        \item[$\rightarrow$] interdependent goods
    \end{itemize}
}
\end{frame}

\begin{frame}{Why Different Production Frontiers?}
\mycomment{
    The theory of comparative advantage assumes different production frontiers, but does not explain why these exist.
    
    Another question is if theses should persist. It may be that focusing on something you can do best locks you into a trade that has no future, and cements an inequality.
}
\mytext{
    \begin{itemize}
        \item `they are better at it'
        \item factor endowment: more of the factor needed present
        \item cheaper labour
        \item higher (labour) productivity
        \item cheaper capital
        \item different preferences
        \item lower standards, laxer laws
        \item path based (`tradition')
        \item should this level out or increase?
    \end{itemize}
}
\end{frame}

